% =====================================================================
%  Podsumowanie i wnioski — co osiągnięto (wnioski odnoszą się do celów
%  z Wstępu!), jakie są ograniczenia i co można zrobić dalej.
%  Rozdział bez numeru, z wpisem w spisie treści.
% =====================================================================
\section*{Podsumowanie i wnioski}
\addcontentsline{toc}{section}{Podsumowanie i wnioski}
\begin{itemize}
	\item Zaprojektowano modułową architekturę biblioteki, umożliwiającą łatwe rozszerzanie
	      przez użytkowników o własne warstwy, optymalizatory, metody inicjalizacji i funkcje straty.
	\item Zaimplementowano od podstaw własny moduł z $n$-wymiarowymi tensorami, obsługującymi
	      mechanizm kroków (\textit{strides}), oraz stosujących podejście typu \textit{zero-copy}
	      (\textit{Copy-on-Write}) -- co pozwalało na uzyskanie wysokiej wydajności oraz efektywnie
	      zarządzanie zasobami. Do operacji mnożenia macierzy skorzystano z zewnętrznej biblioteki.
	\item Stworzono implementacje kluczowych warstw sieci neuronowych -- takich jak warstwy gęste
	      (liniowe), konwolucyjne (Conv2D) z wykorzystaniem własnoręcznie napisanego algorytmu im2col,
	      a także funkcji aktywacji i strat.
	\item W ramach architektury biblioteki opracowano podejście ręcznego różniczkowania -- wymuszającego
	      jawną implementację propagacji wstecznej dla każdego modułu, co pozwala na pełną kontrolę nad przepływem
	      gradientów. Jednocześnie zaimplementowano propagację wsteczną dla warstw wbudowanych.
	\item Stworzono dwa przykładowe programy korzystające z biblioteki -- co pozwoliło na sprawdzenie
	      poprawności jej działania oraz zaprezentowania jej przykładowego użycia.
\end{itemize}
